% RACER-GT Python API reference (English).
\documentclass[12pt,a4paper,oneside]{article}
% English setup. Compiled with XeLaTeX so that both languages share one engine
% and one font family, which keeps the two PDFs typographically comparable.
\usepackage{fontspec}
\setmainfont{Noto Serif}
\setsansfont{Noto Sans}
\setmonofont{Noto Sans Mono}

% Author-year citations. Loaded here rather than in _common because it must
% precede hyperref, and because the Chinese manuscript uses numeric citations
% that natbib's author-year mode would reject. Bibliographies are hand-written
% thebibliography blocks: biblatex/biber is absent from a BasicTeX install, and
% a document set that only builds on a full TeX Live is not reproducible enough
% for this project.
\usepackage[authoryear,round]{natbib}

% Single source of truth for version and date across every RACER-GT document.
% Regenerate nothing by hand: bump here, rebuild, and every PDF agrees.
\newcommand{\racerversion}{1.1.0}
\newcommand{\racerdate}{2026-07-27}
\newcommand{\racerauthor}{Yi-Hao Lai}
\newcommand{\raceraffil}{Department of Finance, Da-Yeh University}
\newcommand{\raceremail}{yhlai@mail.dyu.edu.tw}
\newcommand{\racerrepo}{https://github.com/yhlai0911/RACER-GT}

% Shared package set and mathematical macros for every RACER-GT document,
% in both languages. Language-specific font and theorem-name setup lives in
% _preamble-zh.tex and _preamble-en.tex.
%
% Font policy: the build depends only on the five Noto families below. They are
% installable on Linux CI (fonts-noto-core, fonts-noto-cjk) and on macOS via
% Homebrew casks. Do not introduce a font outside this set without adding it to
% .github/workflows/docs.yml as well, or the PDF build stops being reproducible.

\usepackage[margin=2.35cm,headheight=15pt]{geometry}
\usepackage{amsmath,amssymb,amsthm,mathtools,bm}
\usepackage{booktabs,longtable,tabularx,array,multirow}
\usepackage{graphicx,float,caption,subcaption}
\usepackage{enumitem}
\usepackage{xcolor}
\usepackage{microtype}
\usepackage{setspace}
\usepackage{fancyhdr}
\usepackage{hyperref}
\usepackage[nameinlink,noabbrev]{cleveref}
\usepackage{algorithm}
\usepackage{algpseudocode}
\usepackage{listings}
\usepackage{tikz}
\usetikzlibrary{arrows.meta,positioning,shapes.geometric,fit,calc}

\hypersetup{
  colorlinks=true,
  linkcolor=black,
  citecolor=black,
  urlcolor=blue,
  pdfauthor={\racerauthor}
}

\pagestyle{fancy}
\fancyhf{}
\rhead{v\racerversion}
\cfoot{\thepage}
\setstretch{1.28}
\setlength{\parskip}{0.35em}
\setlist[itemize]{leftmargin=2em,itemsep=0.25em,topsep=0.3em}
\setlist[enumerate]{leftmargin=2.2em,itemsep=0.25em,topsep=0.3em}

% ---------------------------------------------------------------- mathematics
\newcommand{\E}{\mathbb{E}}
\newcommand{\Var}{\operatorname{Var}}
\newcommand{\Cov}{\operatorname{Cov}}
\newcommand{\tr}{\operatorname{tr}}
\newcommand{\diag}{\operatorname{diag}}
\newcommand{\argmin}{\operatorname*{arg\,min}}
\newcommand{\argmax}{\operatorname*{arg\,max}}
\newcommand{\one}{\bm{1}}
\newcommand{\Rset}{\mathbb{R}}
\newcommand{\R}{\mathbb{R}}
\newcommand{\plim}{\operatorname*{plim}}
\newcommand{\RACER}{\textsc{Racer-GT}}
\newcommand{\indic}[1]{\mathbb{I}\!\left\{#1\right\}}
\newcommand{\Sig}{\bm{\Sigma}}
\newcommand{\wvec}{\bm{w}}
\newcommand{\evec}{\bm{e}}
\newcommand{\Zvec}{\bm{Z}}

% Design facets, used identically in both languages so that a reader can move
% between the Chinese and English documents without re-learning notation.
\newcommand{\facetT}{\mathcal{T}}   % historical date
\newcommand{\facetD}{\mathcal{D}}   % collection day
\newcommand{\facetS}{\mathcal{S}}   % collection stream
\newcommand{\facetR}{\mathcal{R}}   % technical replicate

\lstset{
  basicstyle=\ttfamily\footnotesize,
  breaklines=true,
  frame=single,
  rulecolor=\color{gray!40},
  backgroundcolor=\color{gray!5},
  columns=fullflexible,
  keepspaces=true,
  showstringspaces=false
}


\setlength{\parindent}{1.5em}

\newtheorem{assumption}{Assumption}[section]
\newtheorem{proposition}{Proposition}[section]
\newtheorem{corollary}{Corollary}[section]
\newtheorem{remark}{Remark}[section]
\newtheorem{definition}{Definition}[section]
\newtheorem{lemma}{Lemma}[section]

\lhead{RACER-GT: API Reference}

\title{\bfseries RACER-GT Python API Reference\\[0.4em]
\large Public Workflow and Diagnostic Entry Points}
\author{\racerauthor\\\small \raceraffil}
\date{\racerdate\ \textbullet\ version \racerversion}

\begin{document}
\maketitle

\begin{abstract}
\noindent The top-level API is stable. Lower-level functions remain available
for methodological diagnostics but may evolve faster than the pipeline
interface. Where a routine raises rather than returning a degraded result, this
document says so, because those refusals are deliberate design decisions rather
than missing error handling.
\end{abstract}

\tableofcontents
\newpage

\section{Top-level API}

\begin{lstlisting}[language=Python]
from racergt import (
    RacerGTConfig, RacerGTPipeline, PipelineResult,
    DesignWeightedResult, fit_design_weighted_consensus,
    generate_collection_schedule, generate_chunk_windows,
)
\end{lstlisting}

\subsection{\texttt{RacerGTConfig}}

A nested pydantic model holding every locked setting. All sub-models use
\texttt{extra="forbid"}, so an unrecognised YAML key is a validation error
rather than a silently ignored typo.

\begin{lstlisting}[language=Python]
cfg = RacerGTConfig.load_yaml("protocol.yaml")
cfg.save_yaml("protocol.lock.yaml")
digest = cfg.protocol_hash()          # SHA-256 of the canonical serialization
\end{lstlisting}

\texttt{load\_yaml} recomputes the hash and \textbf{raises} if a stored
\texttt{protocol\_hash} disagrees. Hand-editing a lock file therefore fails
loudly. To change a setting, edit the source configuration and regenerate.

\subsection{\texttt{generate\_chunk\_windows(start, end, window\_days, step\_days)}}

Returns the fixed overlapping chunk manifest: \texttt{chunk\_id},
\texttt{window\_start}, \texttt{window\_end}, \texttt{n\_calendar\_days}.

\subsection{\texttt{generate\_collection\_schedule(config, anchor\_date=None)}}

Returns the balanced day $\times$ stream $\times$ replicate schedule with stream
order, planned time slot, deterministic chunk order, fixed historical endpoints,
and the protocol hash.

\begin{remark}
\texttt{collection\_day} in the output is the \emph{ordinal design level}
(0, 1, 2, \dots), not the day offset; \texttt{day\_offset} and
\texttt{planned\_collection\_date} are separate columns.
\end{remark}

\subsection{\texttt{RacerGTPipeline(config).fit(raw, benchmark=None)}}

Runs audit, per-pull overlap calibration, duplicate diagnostics, the
design-based reference estimator, G-studies, covariance-adjusted consensus,
temporal benchmarking when a benchmark is supplied, reliability diagnostics, and
the locked decision tree. Returns a \texttt{PipelineResult}.

Raises \texttt{ValueError} when the raw batch fails the protocol audit, unless
\texttt{stop\_on\_audit\_error=False}.

\subsection{\texttt{PipelineResult}}

\begin{table}[H]
\centering\small
\begin{tabularx}{\textwidth}{lX}
\toprule
Attribute & Contents\\
\midrule
\texttt{final\_series} & final daily index, conditional SE, 95\% interval\\
\texttt{calibrations} & per-pull overlap-graph calibration results\\
\texttt{duplicate\_diagnostics} & exact groups and the near-duplicate graph\\
\texttt{gstudies} & \texttt{level}, \texttt{detection}, \texttt{innovation}\\
\texttt{design\_consensus} & design-based reference estimator\\
\texttt{consensus} & GLS / minimum-variance consensus\\
\texttt{benchmark} & temporal benchmark result, or \texttt{None}\\
\texttt{reliability} & pairwise, day/stream, convergence diagnostics\\
\texttt{decision} & PASS / REVIEW / FAIL and every rule\\
\texttt{audit} & protocol-integrity audit\\
\bottomrule
\end{tabularx}
\end{table}

\texttt{final\_series} returns the benchmark result when one exists and the GLS
consensus otherwise. \texttt{design\_consensus} never feeds it.

\section{Diagnostic entry points}

\subsection{Calibration}

\begin{lstlisting}[language=Python]
from racergt.overlap import OverlapGraphCalibrator, huber_location, OverlapGraphError
\end{lstlisting}

\texttt{OverlapGraphCalibrator.fit} expects exactly one \texttt{series\_id} and
one \texttt{pull\_id}. It raises \texttt{OverlapGraphError} when the overlap
graph is disconnected and \texttt{allow\_disconnected} is false, because a
disconnected graph means the relative scales are only partially identified and
returning a series anyway would conceal that.

\subsection{Duplicates, G-study, consensus}

\begin{lstlisting}[language=Python]
from racergt.duplicates import diagnose_duplicates
from racergt.gstudy import run_gstudy, run_all_gstudies
from racergt.consensus import fit_design_weighted_consensus, fit_gls_consensus
\end{lstlisting}

\texttt{run\_gstudy} raises if the design is not fully crossed or replicate
counts are unequal, rather than estimating components from an unbalanced design.
\texttt{fit\_gls\_consensus} requires at least two pulls and a positive finite
baseline mean for each.

\begin{remark}[Duplicate components are not equivalence classes]
Connected components of the near-duplicate graph are connectivity sets.
Membership does not imply every pair inside exceeds the threshold, because
connectivity is transitive and the pairwise rule is not.
\end{remark}

\subsection{Benchmark and measurement error}

\begin{lstlisting}[language=Python]
from racergt.benchmark import temporal_benchmark
from racergt.eiv import reliability_corrected_ols, simex_ols
\end{lstlisting}

\texttt{reliability\_corrected\_ols} raises when the corrected denominator is
non-positive. That condition means the estimated measurement error accounts for
all observed regressor variation; a number returned there would be meaningless.

\subsection{Simulation and validation}

\begin{lstlisting}[language=Python]
from racergt.simulation import SimulationSettings, simulate_racergt_data
from racergt.validation import run_monte_carlo
\end{lstlisting}

\texttt{run\_monte\_carlo} returns a \texttt{ValidationResult} holding
\texttt{metrics} (one row per replication, estimator and information set),
\texttt{summary} (means plus \texttt{sd\_rmse} and \texttt{mcse\_rmse}) and
\texttt{paired\_tests}. Each estimator is evaluated twice, once on the calibrated
pulls alone and once after applying the \emph{same} temporal benchmark, because
comparing a benchmarked \RACER{} against unbenchmarked baselines confounds the
estimator with the extra information. The single-pull row averages over all
pulls rather than taking a fixed column. \texttt{scripts/make\_figures.py}
redraws the manuscript figures from the same CSVs.

\section{File I/O}

\texttt{racergt.io.read\_table} reads CSV, Parquet, XLS/XLSX;
\texttt{write\_table} writes CSV, Parquet, XLSX, and Stata 118 \texttt{.dta}.

\section{Command line}

\begin{table}[H]
\centering\small
\begin{tabularx}{\textwidth}{lXl}
\toprule
Command & Purpose & Exit\\
\midrule
\texttt{racergt design} & lock protocol, schedule, chunk manifest & 0\\
\texttt{racergt audit} & check a batch against the locked protocol & 2 on failure\\
\texttt{racergt run} & run the full pipeline & 3 on FAIL\\
\texttt{racergt simulate} & controlled dataset, optionally run pipeline & 0\\
\texttt{racergt export-stata} & CSV to Stata 118 & 0\\
\bottomrule
\end{tabularx}
\end{table}

\end{document}
